\begin{abstract}
The base of a supernumerary robotic limb is a person. The aim of this project
was to build a teleoperation system for two seven-degree-of-freedom arms
carried on a worn backpack frame and commanded by a \emph{second} person, give
it a safety layer appropriate to a manipulator standing beside somebody who is
not commanding it, and characterise what it can and cannot do.

A body-worn master supplies position from whichever sensor observes each
component most directly, and classifies all fourteen of its channels so that a
failing sensor degrades the system announcedly rather than silently. Sixteen
simulations, run by one command, and four physical cameras measured the
result.

Six of fourteen master channels met the health criterion, and the failing ones
injected \SI{54.5}{\milli\metre} of commanded motion nobody asked for.
Vertical and fore-and-aft command track correctly; lateral does not, because
directions intended to be \SI{90}{\degree} apart measure
\SI{137.7}{\degree} apart, and no axis mapping makes a non-orthogonal triad
orthogonal. The two arms share no reachable point, which blocks inter-arm
handover for geometric reasons no control change addresses. The task set solves \num{13927} inverse-kinematics calls with no failures and
no breach of the \SI{150}{\milli\metre} clearance floor around the wearer;
replacing the motion generator removed \SI{51.3}{\milli\metre} of unrequested
excursion; and measuring object positions rather than declaring them takes the
autonomy layer from confidently wrong on half of trials to never.

Nothing was validated on a loaded physical arm, and the two-person study the
architecture was built for was designed in full but could not be run.
\end{abstract}

\renewcommand{\abstractname}{Acknowledgements}
\begin{abstract}
I thank my supervisor, Dr.\ Dandan Zhang, for the direction of this project
and for consistently pushing the work toward what could be measured rather
than what could be claimed. Much of what is useful in this thesis exists
because a result was questioned before it was written down.

I thank my co-supervisor, Dr.\ Etienne Burdet, for the framing that changed
the shape of the project. Treating the wearer and the operator as two people
with different stakes, rather than as one role, is what turned a control
problem into a question worth asking, and it is the reason the workspace and
safety chapters exist in the form they do.

I am grateful to the Human Robotics Group for access to the Dr.\ Octopus
platform, and to everyone who kept two Kinova arms and a backpack frame
available and working.

This project spent a great deal of its time discovering that the instrument
was wrong rather than the system. I am thankful to everyone who responded to a
surprising number with scepticism instead of enthusiasm, which is the only
reason several of them are not in this document as findings.

Finally, I thank my family and friends for their support over the course of
the year.
\end{abstract}
